\section{Introduction}

A knowledge graph that more than one writer can edit will diverge, and something
has to put it back together. The obvious move is to reuse the tooling that
already solves this for source code: keep the graph in a file, keep the file in
git, and let a line merge reconcile two versions.

That move fails in a specific and instructive way. A line merge understands a
document as an ordered sequence of lines, and a graph is neither ordered nor a
sequence. Two writers who add unrelated facts about the same entity produce
adjacent edits and a textual conflict. Two writers who re-serialise the same
unchanged graph produce a diff of everything. Meanwhile a writer who sets a
second value on a property the ontology says holds exactly one --- the change
that genuinely cannot be reconciled without a decision --- produces a clean,
conflict-free merge, because the two values live on different lines and nothing
in the file says they compete.

The failure is not that line merge is bad at graphs. It is that conflict is a
\emph{semantic} property of the data model, and a line merge is being asked to
infer it from a syntactic one. Our position is that an RDF merge should take
conflict from the place the data model already defines it: the schema.

\paragraph{This paper.} We describe an operator with two halves. Triples are
canonicalised and merged as sets, so serialisation order cannot manufacture a
conflict. A SHACL shapes graph then decides which slots the set algebra is
allowed to settle on its own: where a shape declares \texttt{sh:maxCount 1} and
the two sides carry different values, the operator refuses and records a
decision; everywhere else, it unions. Because the shapes are already present ---
they are what validates writes into the store --- the merge policy is not a
second configuration to keep in step with the ontology. It \emph{is} the
ontology.

\paragraph{What we are not claiming.} Three-way merge over canonicalised RDF is
not new, and we searched before writing. Quit Store~\cite{arndt2019quit} places
RDF under git with merge strategies, including a Context Merge that identifies
conflicts by node overlap and refers them to a person; TerminusDB offers diff and patch over a versioned graph; RDF
Patch and LD Patch standardise change description; RDFC-1.0 gives canonical
labelling for graphs with blank nodes, which we build on rather than beside. To
our knowledge --- and scoped to that search --- what has not been reported is an
evaluation of \emph{SHACL-derived} conflict semantics against those neighbours
in a deployed system. That evaluation is the contribution, and the framing of
the operator as an extension of Context Merge rather than a replacement for it
is deliberate.

\paragraph{Contributions.}
\begin{enumerate}[leftmargin=*,itemsep=2pt]
  \item A three-way merge operator over canonical triples whose conflict
        predicate is derived from SHACL cardinality, with post-merge validation
        against the same shapes, implemented in a production store
        (\S\ref{sec:operator}, \S\ref{sec:impl}).
  \item A divergence benchmark with an oracle recomputable by a reader from
        published data alone, scoring eight strategies under one accounting
        contract, and a result that contradicts one of our own hypotheses
        (\S\ref{sec:synthetic}).
  \item A replay of recorded multi-agent divergence from a production graph
        against the shipped operator, which reproduces the synthetic arm's
        blind spot on real data and quantifies the benefit on a class that
        production surfaced nothing of (\S\ref{sec:replay}).
  \item A stated limit --- alias divergence is invisible to triple-level merge
        --- reported as a measured false-negative column rather than excluded
        from the benchmark (\S\ref{sec:limits}).
\end{enumerate}
